% =============================================================================
% Section 8: Expected Results
% =============================================================================

\section{Expected Results}
\label{sec:expected}

% AGENT: Fill expected limit value and Brazil band plot from Wave 4
% This section presents the full set of expected results before unblinding.

\subsection{Expected Upper Limit}
\label{sec:expected:limit}

% AGENT: Fill expected limit from Fitter.expected_limit() output.
% Include central value and +/-1,2 sigma bands.

The expected 95\% CL$_\text{s}$ upper limit on the signal strength parameter
$\mu$ is determined using the asymptotic approximation with an Asimov dataset.
The results are summarized in Table~\ref{tab:expected_limit}.

\begin{table}[htbp]
  \centering
  \caption{Expected 95\% CL$_\text{s}$ upper limit on the signal strength
           $\mu$ from the Asimov fit.}
  \label{tab:expected_limit}
  \begin{tabular}{lr}
    \toprule
    Quantity & Value \\
    \midrule
    % AGENT: Fill from Fitter.expected_limit() output
    $-2\sigma$ expected & [value] \\
    $-1\sigma$ expected & [value] \\
    Median expected     & [value] \\
    $+1\sigma$ expected & [value] \\
    $+2\sigma$ expected & [value] \\
    \bottomrule
  \end{tabular}
\end{table}

\subsection{Fit Diagnostics}
\label{sec:expected:diagnostics}

% AGENT: Fill fit diagnostic results from Diagnostics.run_all_diagnostics()
% on the Asimov dataset. Include goodness-of-fit, NP pulls, and constraints.

The quality of the statistical model is assessed through several diagnostic
tests performed on the Asimov dataset.

\subsubsection{Nuisance Parameter Pulls}

% AGENT: Include NP pull plot from Diagnostics output.
% Verify all pulls are within +/-2 sigma of nominal.

\begin{figure}[htbp]
  \centering
  % AGENT: Replace with actual pull plot path from Wave 4 diagnostics
  % \includegraphics[width=0.8\textwidth]{figures/np_pulls_asimov.pdf}
  \caption{Post-fit nuisance parameter pulls from the Asimov fit.
           All parameters are expected to be consistent with their
           pre-fit values.}
  \label{fig:np_pulls}
\end{figure}

\subsubsection{Nuisance Parameter Ranking}

% AGENT: Include NP ranking plot showing impact on mu.
% Top-ranked NPs indicate dominant systematic effects.

\begin{figure}[htbp]
  \centering
  % AGENT: Replace with actual ranking plot path from Wave 4 diagnostics
  % \includegraphics[width=0.8\textwidth]{figures/np_ranking_asimov.pdf}
  \caption{Ranking of nuisance parameters by their impact on the signal
           strength $\mu$, evaluated from the Asimov fit.}
  \label{fig:np_ranking}
\end{figure}

\subsubsection{Correlation Matrix}

% AGENT: Include correlation matrix plot for the most correlated NPs.

\begin{figure}[htbp]
  \centering
  % AGENT: Replace with actual correlation matrix path from Wave 4 diagnostics
  % \includegraphics[width=0.7\textwidth]{figures/correlation_matrix_asimov.pdf}
  \caption{Correlation matrix of the most correlated nuisance parameters
           from the Asimov fit.}
  \label{fig:correlation}
\end{figure}

\subsubsection{Likelihood Scan}

% AGENT: Include likelihood scan plot for the POI (mu).

\begin{figure}[htbp]
  \centering
  % AGENT: Replace with actual likelihood scan path from Wave 4 diagnostics
  % \includegraphics[width=0.6\textwidth]{figures/likelihood_scan_asimov.pdf}
  \caption{Profile likelihood ratio scan as a function of the signal
           strength $\mu$ for the Asimov dataset.}
  \label{fig:likelihood_scan}
\end{figure}

\subsection{Sensitivity Optimization}
\label{sec:expected:optimization}

% AGENT: Fill sensitivity optimization summary from SensitivityOptimizer results.
% Document any selection or categorization adjustments made to improve sensitivity.

The analysis sensitivity was optimized by [optimization procedure description].
The final expected limit of $\mu < $ [value] at 95\% CL represents an
improvement of [factor] compared to [baseline/previous configuration].
